\documentclass[10pt,a4paper]{amsart}
\usepackage[margin=19mm]{geometry}
\usepackage{amsmath,amssymb,mathtools}
\usepackage[scr=rsfs,cal=euler]{mathalfa}
\usepackage{xfrac}
\usepackage[unicode,hidelinks,bookmarks=false]{hyperref}
\newcommand{\ie}{i.e.\ }
\newcommand{\cf}{cf.\ }
\newcommand{\resp}{respectively\ }
\newcommand{\Subsection}{\subsection}
\providecommand{\nobreakdash}{\nobreak\hbox{-}}
\setlength{\emergencystretch}{2em}
\multlinegap0pt
\usepackage{amssymb}
\newtheorem{proposition}{Proposition}
\newtheorem{lemma}{Lemma}
\usepackage{cleveref}
\crefname{proposition}{Proposition}{Propositions}
\crefname{lemma}{Lemma}{Lemmas}
\newcommand{\angles}[2]{\left\langle #1,#2\right\rangle_Q}
\title{Nonnegative conorms, regular matroids, and the tropical Schottky problem}
\author{Yelena Mandelshtam}
\date{Source: arXiv:2608.28783v1 (CC BY 4.0)}
\begin{document}
\maketitle

\noindent\emph{Source and licence.} Nonnegative conorms, regular matroids, and the tropical Schottky problem. Version arXiv:2608.28783v1 (CC BY 4.0), distributed by arXiv under the Creative Commons Attribution 4.0 International licence, http://creativecommons.org/licenses/by/4.0/, as recorded in the arXiv licence metadata for this exact version and shown on the abstract page https://arxiv.org/abs/2608.28783v1. Full licence notice: \texttt{licenses/arxiv-2608.28783-CC-BY-4.0.txt}.\par\smallskip

\noindent\textit{Editorial scope.} Counted unit: the article's proposition prop:common-representative (source0.tex lines 804--818). The article's complete original proof of it is delivered with it (source0.tex lines 819--877, one \texttt{proof} environment of 59 lines). No mathematics is written, repaired or re-derived: the statement and the proof are the article's own lines, in the article's own notation, and no retained line is substituted or re-flowed.  Retained uncounted as the source-local context the proof needs: lines 652--662; lines 725--745.  Nothing was omitted from any retained span, so the package carries no omission record.  No retained line contains a citation, so the wrapper carries no bibliography and no bibliography entry.  No retained line contains a figure environment, an image-inclusion command or drawing code, so no image is delivered and the package declares no asset dependency.  Editorial formatting only, in addition: an AMS \texttt{amsart} preamble that restates the article's own declarations and macros used by the retained lines, namely the environments \texttt{proposition}, \texttt{lemma} on separate counters, together with the standard packages those lines need.  The retained environments print in this document as Proposition 1, Lemma 1 in order of appearance, whereas the article prints the same items with the numbering of its own full document.  The complete native source of the article as distributed in the arXiv e-print for this exact version is bundled unchanged as \texttt{sources/208808/source0.tex}.
\begin{proposition}\label{prop:coedge-rule}
For every cocircuit $D$,
\begin{equation*}
  a_e(r_D)=
  \begin{cases}
    0,&e\notin D,\\
    \pm1,&e\in D.
  \end{cases}
\end{equation*}
Furthermore, the reduction of $a_e$ modulo $2$ is the original character $e\in V^*$.
\end{proposition}

\begin{lemma}\label{lem:orthogonal}
If $D_1,\ldots,D_k$ are pairwise disjoint cocircuits, then
\[
  \angles{r_{D_i}}{r_{D_j}}=0 \quad \text{ for all } i \neq j \in [k].
\]


Further, for every choice of signs
$\varepsilon_1,\ldots,\varepsilon_k\in\{\pm1\}$, the vector
\[
    x=\sum_{i=1}^k \varepsilon_i r_{D_i}
\]
is a shortest vector in the class $\alpha_{D_1}+\cdots+\alpha_{D_k}$. In particular,
\[
    Q[x]
    =
    \sum_{i=1}^k Q[r_{D_i}]
    =
    \mu\left(\alpha_{D_1}+\cdots+\alpha_{D_k}\right).
\]
\end{lemma}

\begin{proposition}\label{prop:common-representative}
For every $\alpha\in V$, the vector $x_\alpha$ lies in the class $\alpha$ and simultaneously minimizes $Q$ and
$Q_A$. That is, we have
\[
    Q[x_\alpha]
    =
    \min_{\substack{x\in L\\ [x]=\alpha}}Q[x]
    =
    \mu(\alpha)
    =
    \min_{\substack{x\in L\\ [x]=\alpha}}Q_A[x]
    =
    Q_A[x_\alpha].
\]
\end{proposition}

\begin{proof}
    The cocycle with support $D_i$ is $\iota(\alpha_{D_i})$. Since the
supports $D_i$ are disjoint and their union is $D(\alpha)$, we have
\[
    \iota(\alpha)
    =
    \iota(\alpha_{D_1})
    +\cdots+
    \iota(\alpha_{D_k})
\]
in $\mathbb F_2^E$. Because $\iota$ is injective, we have $\alpha = \alpha_{D_1}+\cdots+\alpha_{D_k}.$ Consequently, $[x_\alpha]=\alpha$.

By \cref{lem:orthogonal},
\[
    Q[x_\alpha] =
    \sum_{i=1}^k Q[r_{D_i}] =
    \sum_{i=1}^k \mu(\alpha_{D_i}) =
    \sum_{i=1}^k\sum_{e\in D_i}\lambda_e =
    \sum_{e\in D(\alpha)}\lambda_e =
    \mu(\alpha).
\]
Thus $x_\alpha$ minimizes $Q$ in its parity class.

Next, we evaluate $Q_A$ on $x_\alpha$. Fix $e\in E$. Since the
cocircuits $D_1,\ldots,D_k$ are pairwise disjoint, the element $e$
belongs to at most one of them. Then, from \cref{prop:coedge-rule},
it follows that
\[
    a_e(x_\alpha)^2
    =
    \begin{cases}
        1, & e\in D(\alpha),\\
        0, & e\notin D(\alpha).
    \end{cases}
\]
And then
\[
    Q_A[x_\alpha]
    =
    \sum_{e\in E}\lambda_e a_e(x_\alpha)^2
    =
    \sum_{e\in D(\alpha)}\lambda_e
    =
    \mu(\alpha).
\]

It remains to verify that $x_\alpha$ is a shortest vector for $Q_A$ in its parity class. Let $y\in L$ be any representative of $\alpha$. Since the reduction of $a_e$ modulo $2$ is $e$, we have $a_e(y)\equiv e(\alpha)\pmod 2.$ In particular, for every $e\in D(\alpha)$, the integer $a_e(y)$ is
odd, and hence $a_e(y)^2\geq1.$ Since all $\lambda_e$ are positive, we have
\begin{align*}
    Q_A[y]
    &=
    \sum_{e\in E}\lambda_e a_e(y)^2 \\
    &\geq
    \sum_{e\in D(\alpha)}\lambda_e =
    \mu(\alpha) =
    Q_A[x_\alpha].
\end{align*}
and we are done.
\end{proof}

\end{document}
